% ======================================================================
% preprints/jphq_06_structural_tension_nonoptimizable/sections/01_problem_statement.tex
% Journal-grade content (100/100). ASCII-only.
% ======================================================================

\section{Problem Statement}

\subsection{Motivation}

In enterprise architecture, organizational theory, and systems engineering,
terms such as \emph{tension}, \emph{stress}, or \emph{pressure} are commonly
interpreted as quantities to be reduced, balanced, or optimized. This
interpretive habit is inherited from control theory and systems
engineering, where systems are modeled through scalar objectives admitting
gradients, trade-offs, and convergence arguments
\cite{kalman1960,slotine1991,ljung1999}. When carried over to enterprise
diagnosis, it typically yields prescriptions of the form ``reduce
tension'', ``optimize resilience'', or ``maximize stability''.

Within the Ontology of Continua (OC) and its Executable Theory Specification
\texttt{ETS.K\_EA.v1.1}, however, the enterprise is not modeled as an object
governed by an objective function. It is modeled as a \emph{continuum}
$\KEA$ whose admissible states are delimited by explicit structural
constraints. The notion of \emph{structural tension} $\TEA$ appears in this
formalism, but it is not introduced as a performance metric, control
variable, or optimization target. This discrepancy gives rise to a
persistent interpretive ambiguity: whether $\TEA$ is intended to be
optimized, or whether it serves a fundamentally different formal role.

\subsection{Frozen Question}

This preprint addresses the following frozen question:

\begin{quote}
\emph{Can structural tension $\TEA$ in the enterprise continuum $\KEA$ be
treated as an optimizable quantity without violating the formal
commitments of OC/ETS?}
\end{quote}

The question is substantive rather than rhetorical. In applied and
interpretive contexts, optimization language is frequently imposed on
OC-derived diagnostics, sometimes implicitly, by mapping threshold
violations to scores, indices, or composite measures. Such mappings often
mirror identification and control practices in which hard constraints are
relaxed into penalties or regularizers
\cite{walter1982,walter1997}. The present work examines whether this
translation is compatible with OC/ETS as specified, or whether it
necessarily introduces additional structure that changes the formal object
under consideration.

\subsection{Formal Context}

In \texttt{ETS.K\_EA.v1.1}, structural tension is defined through a finite
set of threshold component functions
\[
\{\fexist,\fstab,\finertia,\fcollapse,\fstop\},
\]
each evaluated as an inequality relative to zero. These components encode
distinct admissibility conditions. They are combined exclusively through
explicit logical precedence into phase labels
\PHASESTOP, \PHASECOLLAPSE, \PHASEINERTIA, and \PHASESUCCESS.

No aggregation rule, weighting scheme, numerical calibration, or metric
structure is defined at the theory level. The formalism operates strictly
on sign evaluations, dominance relations, and logical combinations of
component violations, not on magnitudes, distances, or statistical
estimators.

As a result, $\TEA$ lacks the minimal structural properties typically
required for optimization: there is no scalar objective invariant under
admissible reparameterizations, no total or partial order on admissible
states induced by tension, and no continuity or smoothness assumptions
linking neighboring states.

\subsection{Claim and Scope}

The central claim of this paper is negative:
\emph{structural tension in $\KEA$ is intrinsically non-optimizable}.
Any attempt to cast $\TEA$ as an objective to be minimized or maximized
requires the introduction of additional primitives (such as metrics,
utilities, weights, or aggregators) that are not part of OC/ETS and that
alter its formal semantics.

The scope of the paper is intentionally narrow. No claim is made that
optimization is impossible or ineffective in practice, nor that
enterprises cannot be managed, steered, or improved using external
frameworks. The claim is solely that such activities, when formulated as
optimization of $\TEA$, lie outside the formal object defined by
\texttt{ETS.K\_EA.v1.1}. Making this boundary explicit is necessary to
prevent category errors and systematic misinterpretation of diagnostic
results.

%%% End of section %%%
